% ===========================================================================
%  北大《高等代数》第五版 · 第四章 矩阵 §4 矩阵的逆 —— 深化提取
%  （服务于蓝本第 2 章「余子式和代数余子式的计算」，考纲「矩阵」：
%    理解伴随矩阵的概念，会用伴随矩阵求逆矩阵）
%  源：PDF p131–p132 = 印刷页 119–120
%      实测：p131、p132 页眉均为「§4 矩阵的逆」；p131 页脚 = 119，
%      p132 页脚 = 120，本段偏移 = PDF 页 − 12。
%  原则：只取定义 / 定理 / 推论 / 关键证明思路；例题、习题、纯计算演示全部跳过。
% ===========================================================================


% ---------------------------------------------------------------------------
% 【深化】服务考点：矩阵的逆（概念与唯一性）
% ---------------------------------------------------------------------------
\begin{fdeep}{逆矩阵的定义与唯一性}
\fsec{§4\quad 矩阵的逆}

首先我们指出，由于矩阵的乘法规则，只有方阵才能满足 $AB=BA=E$（读者自己证明）.其次，对于任意的矩阵 $A$，适合等式 $AB=BA=E$ 的矩阵 $B$ 是唯一的（如果有的话）.事实上，假设 $B_1,B_2$ 是两个适合它的矩阵，就有
\fml{B_1=B_1E=B_1(AB_2)=(B_1A)B_2=EB_2=B_2.}

\textbf{定义 8}\quad 如果矩阵 $B$ 适合 $AB=BA=E$，那么 $B$ 就称为 $A$ 的逆矩阵，记作 $A^{-1}$.下面要解决的问题是：在什么条件下矩阵 $A$ 是可逆的？如果 $A$ 可逆，怎样求 $A^{-1}$？

【反哺点】服务考纲「矩阵」中"理解逆矩阵的概念、掌握逆矩阵的性质"：唯一性的证明只用到"$E$ 是乘法单位元 $+$ 结合律"，是"证唯一性"这类小题的标准四步；同时点明逆矩阵只对方阵才有意义，可避免对非方阵误求逆。
\end{fdeep}


% ---------------------------------------------------------------------------
% 【深化】服务考点：伴随矩阵的定义（考纲「矩阵」明确要求"理解伴随矩阵的概念"）
% 反哺点：$A^*$ 的元素按"代数余子式再转置"摆放，是伴随矩阵最高频的失分点
% ---------------------------------------------------------------------------
\begin{fdeep}{伴随矩阵的定义：代数余子式的转置排布（定义 9）}
\textbf{定义 9}\quad 设 $A_{ij}$ 是矩阵
\fml{A=\begin{pmatrix}
a_{11} & a_{12} & \cdots & a_{1n}\\
a_{21} & a_{22} & \cdots & a_{2n}\\
\vdots & \vdots & & \vdots\\
a_{n1} & a_{n2} & \cdots & a_{nn}
\end{pmatrix}}
中元素 $a_{ij}$ 的代数余子式，矩阵
\fml{A^{*}=\begin{pmatrix}
A_{11} & A_{21} & \cdots & A_{n1}\\
A_{12} & A_{22} & \cdots & A_{n2}\\
\vdots & \vdots & & \vdots\\
A_{1n} & A_{2n} & \cdots & A_{nn}
\end{pmatrix}}
称为 $A$ 的伴随矩阵.

【反哺点】服务考纲「矩阵」"理解伴随矩阵的概念，会用伴随矩阵求逆矩阵"：$A^{*}$ 的第 $i$ 行第 $j$ 列元素是 $A_{ji}$，即"代数余子式矩阵的转置".若误记成 $A_{ij}$ 原位摆放，二阶以上全部算错——这是伴随矩阵题的第一失分点，务必与"转置"绑定记忆。
\end{fdeep}


% ---------------------------------------------------------------------------
% 【深化】服务考点：余子式 / 代数余子式的计算（蓝本第 2 章）→ 矩阵
% 反哺点：$AA^{*}=A^{*}A=|A|E$ 的证明思路，正是"一行元素乘另一行代数余子式之和为 0"
% ---------------------------------------------------------------------------
\begin{fdeep}{伴随矩阵基本恒等式 $AA^{*}=A^{*}A=|A|E$}
由行列式按一行（列）展开的公式立即得出
\fml{AA^{*}=A^{*}A=\begin{pmatrix}
d & 0 & \cdots & 0\\
0 & d & \cdots & 0\\
\vdots & \vdots & & \vdots\\
0 & 0 & \cdots & d
\end{pmatrix}=dE,\qquad (2)}
其中 $d=|A|$.

证明思路（按元素展开）：$(AA^{*})$ 的 $(i,j)$ 元素是 $A$ 的第 $i$ 行与 $A^{*}$ 的第 $j$ 列对应相乘之和，而 $A^{*}$ 的第 $j$ 列恰是 $A$ 第 $j$ 行元素的代数余子式 $A_{j1},A_{j2},\cdots,A_{jn}$，于是
\fml{(AA^{*})_{ij}=\sum_{k=1}^{n}a_{ik}A_{jk}=
\begin{cases}
d, & i=j,\\[2pt]
0, & i\ne j.
\end{cases}}
前者是行列式按第 $i$ 行的展开式（得 $d=|A|$），后者是"一行的元素乘另一行对应元素的代数余子式之和为零"（得 $0$）.对 $A^{*}A$ 按列展开同理.

【反哺点】服务考纲「矩阵」与蓝本第 2 章"余子式和代数余子式的计算"：这条恒等式是把"代数余子式"从一个行列式技巧升级为矩阵工具的唯一桥梁——对角元来自展开定理，非对角元来自"异行代数余子式之和为 0".记住这个"一正一零"的来源，$AA^{*}=|A|E$ 就不用背，也能自己现场推。
\end{fdeep}


% ---------------------------------------------------------------------------
% 【深化】服务考点：可逆的充要条件 + 用伴随矩阵求逆（考纲「矩阵」核心要求）
% ---------------------------------------------------------------------------
\begin{fdeep}{矩阵可逆的充要条件与逆矩阵公式（定理 3）}
如果 $d=|A|\ne 0$，那么由 (2) 得
\fml{A\left(\frac{1}{d}A^{*}\right)=\left(\frac{1}{d}A^{*}\right)A=E.\qquad (3)}

\textbf{定理 3}\quad 矩阵 $A$ 是可逆的充分必要条件是 $A$ 非退化，而
\fml{A^{-1}=\frac{1}{d}A^{*}\qquad (d=|A|\ne 0).}

证明思路（充分性）：当 $d=|A|\ne 0$ 时，由 (3) 可知 $A$ 可逆，且
\fml{A^{-1}=\frac{1}{d}A^{*}.\qquad (4)}

证明思路（必要性）：反过来，如果 $A$ 可逆，那么有 $A^{-1}$ 使 $AA^{-1}=E$.两边取行列式，得
\fml{|A|\,|A^{-1}|=|E|=1,\qquad (5)}
因而 $|A|\ne 0$，即 $A$ 非退化.

【反哺点】服务考纲「矩阵」"会用伴随矩阵求逆矩阵"：$A^{-1}=\dfrac{1}{|A|}A^{*}$ 就是考纲原话的公式形态，二阶情形 $\begin{pmatrix}a&b\\c&d\end{pmatrix}^{-1}=\dfrac{1}{ad-bc}\begin{pmatrix}d&-b\\-c&a\end{pmatrix}$ 由它直接得到.证明演示了两条通用手法——"先凑出 $AX=E$ 即可断言可逆"与"等式两边取行列式".而"$|A|\ne0\Leftrightarrow A$ 可逆"是贯穿方程组、秩、特征值讨论的总开关.
\end{fdeep}


% ---------------------------------------------------------------------------
% 【深化】服务考点：逆矩阵的判定与行列式性质
% ---------------------------------------------------------------------------
\begin{fdeep}{只验一边即可判定互逆，以及 $|A^{-1}|=|A|^{-1}$}
根据定理 3 容易看出，对于 $n$ 阶方阵 $A,B$，如果
\fml{AB=E,}
那么 $A,B$ 就都是可逆的并且它们互为逆矩阵.

由 (5) 可以看出，如果 $d=|A|\ne 0$，那么
\fml{|A^{-1}|=d^{-1}.}

（另注：定理 3 不但给出了矩阵可逆的条件，同时也给出了求逆矩阵的公式 (4).按这个公式来求逆矩阵，计算量一般是很大的.）

【反哺点】服务考纲「矩阵」：判定 $A,B$ 互逆只需验证一边 $AB=E$，不必再验 $BA=E$——选项题与证明题里省掉一半工作量；$|A^{-1}|=|A|^{-1}$ 常与 $|A^{*}|=|A|^{n-1}$、$|kA|=k^{n}|A|$ 串联出现，是行列式数值计算的固定链条；末句提示伴随公式适合理论推导而非大阶数手算（大阶数应改用初等行变换法），可避免考场上硬算伴随矩阵。
\end{fdeep}


% ---------------------------------------------------------------------------
% 【深化】服务考点：逆矩阵的运算性质（转置、乘积）
% ---------------------------------------------------------------------------
\begin{fdeep}{逆矩阵的运算性质：转置与乘积（推论）}
\textbf{推论}\quad 如果矩阵 $A,B$ 可逆，那么 $A^{\mathrm{T}}$ 与 $AB$ 也可逆，且
\fml{(A^{\mathrm{T}})^{-1}=(A^{-1})^{\mathrm{T}},\qquad (AB)^{-1}=B^{-1}A^{-1}.}

证明思路：由定理即得推论的前一半，现在来证后一半.由
\fml{AA^{-1}=A^{-1}A=E}
两边取转置，有
\fml{(A^{-1})^{\mathrm{T}}A^{\mathrm{T}}=A^{\mathrm{T}}(A^{-1})^{\mathrm{T}}=E^{\mathrm{T}}=E,}
因此 $(A^{\mathrm{T}})^{-1}=(A^{-1})^{\mathrm{T}}$.由
\fml{(AB)(B^{-1}A^{-1})=(B^{-1}A^{-1})(AB)=E}
即得 $(AB)^{-1}=B^{-1}A^{-1}$.

【反哺点】服务考纲「矩阵」"掌握逆矩阵的性质"：证明给出两条通用手段——"等式两边取转置"与"直接验证乘积等于 $E$"（配合上一条的"只验一边"）.$(AB)^{-1}=B^{-1}A^{-1}$ 的顺序颠倒与 $(AB)^{\mathrm{T}}=B^{\mathrm{T}}A^{\mathrm{T}}$ 完全同构，是求分块矩阵逆、解矩阵方程 $AXB=C$（得 $X=A^{-1}CB^{-1}$）时的关键，绝不能写成 $A^{-1}B^{-1}$.
\end{fdeep}


% ---------------------------------------------------------------------------
% 【深化】服务考点：线性方程组（克拉默法则的矩阵推导）
% ---------------------------------------------------------------------------
\begin{fdeep}{用逆矩阵推导克拉默法则}
利用矩阵的逆，可以给出克拉默法则的另一种推导法.线性方程组
\fml{\begin{cases}
a_{11}x_1+a_{12}x_2+\cdots+a_{1n}x_n=b_1,\\
a_{21}x_1+a_{22}x_2+\cdots+a_{2n}x_n=b_2,\\
\cdots\cdots\cdots\cdots\\
a_{n1}x_1+a_{n2}x_2+\cdots+a_{nn}x_n=b_n
\end{cases}}
可以写成
\fml{AX=B.\qquad (6)}
如果 $|A|\ne 0$，那么 $A$ 可逆.用
\fml{X=A^{-1}B}
代入 (6)，得恒等式 $A(A^{-1}B)=B$，这就是说 $A^{-1}B$ 是一个解.

如果 $X=C$ 是 (6) 的一个解，那么由 $AC=B$ 得
\fml{A^{-1}(AC)=A^{-1}B,}
即
\fml{C=A^{-1}B.}
这就是说，解 $X=A^{-1}B$ 是唯一的.用 $A^{-1}$ 的公式 (4) 代入，乘出来就是克拉默法则中给出的公式.

【反哺点】服务考纲「线性方程组」"会用克拉默法则"：把"解方程组"统一成"求逆矩阵"，正是考纲「矩阵」与「线性方程组」两节的接口.这段还给出"存在性（代回验证 $A(A^{-1}B)=B$）$+$ 唯一性（两边左乘 $A^{-1}$）"的完整证明模板，凡是问"证明某矩阵表达式是唯一解"的题都可照此两步走.
\end{fdeep}
